\section{관련 연구}

\paragraph{대화형 AI 에이전트를 위한 벤치마크.}
언어 에이전트(language agent)~\cite{yao2022webshop, zhou2023webarena, jimenez2023swe, liu2023agentbench, ruan2023identifying}, LLM 도구 사용~\cite{berkeley-function-calling-leaderboard,qin2024toolllm,huang2023metatool}, 과제 지향 대화(task-oriented dialog)~\cite{chen2021action, budzianowski2018multiwoz, andreas2020task,schatzmann2007evaluating,gur2018hidden,he2018decoupling,Hu_2023}에 관한 긴 연구 흐름을 따라, \taubench~\cite{yao2024tau}는 도메인 규칙을 준수하면서 고객 서비스 워크플로와 같은 다중 턴 과제 지향 대화에서 언어 에이전트의 신뢰성을 측정하기 위해 최근 도입된 벤치마크다. \taubench의 각 과제는 사용자 시뮬레이터와 언어 에이전트 사이의 실시간 대화를 인스턴스화하며, 과제는 \retail 및 \airline 두 도메인에 걸쳐 있다. 신뢰성을 정량화하기 위해 이 논문은 \passk 지표, 즉 \textit{k}개의 독립 실행 중 성공한 비율을 도입한다.

\taubench의 여러 후속 연구는 기본 설정의 변형을 탐구했다. FlowBench~\cite{xiao2024flowbench}는 자연어, python 유사 의사코드 또는 mermaid 플로차트를 사용해 명시적 워크플로 지식을 프롬프트에 주입함으로써 도구 사용 에이전트의 계획 단계를 분리한다. IntellAgent~\cite{levi2025intellagent}는 도메인 규칙과 그 동시 발생 통계를 인코딩한 구조화된 정책 그래프로부터 합성 테스트 스위트를 프로그램적으로 구축하는 평가 파이프라인을 제공한다. IntellAgent는 \taubench를 외부 골드 스탠더드로 명시적으로 사용하며, 두 점수 분포 사이의 높은 Spearman 상관을 보고하고 빠른 합성 프록시 과제로 기능한다. APIGen-MT~\cite{prabhakar2025apigen}는 \taubench를 위해 도구 호출 에이전트를 파인튜닝(fine-tuning)하는 아이디어를 탐구한다. 이들은 서로 의존하는 도구 호출의 시퀀스인 대화 청사진을 만들고, 각 청사진에 기반해 대화 추적을 시뮬레이션함으로써 데이터를 생성한다. ToolSandbox~\cite{lu2024toolsandbox}는 에이전트 진행 상황을 더 세분화된 방식으로 평가하기 위해 상태를 갖는 도구를 만드는 데 초점을 맞춘다.

우리 연구는 \taubench 패러다임을 확장하고, 사용자와 에이전트가 모두 공유 세계 위에서 (도구 호출을 통해) 상태를 바꾸는 능력을 가질 수 있도록 일반화한다. 결과에서 보이듯이 이는 대화형 에이전트를 시험하기 위한 더 복잡한 도메인을 구축할 수 있게 하며, 개선 가능한 에이전트 실패 지점의 세분화된 분석 기회도 제공한다.

\paragraph{대화형 에이전트를 위한 사용자 시뮬레이션.}
사용자 시뮬레이션의 신뢰성은 \taubench~\cite{yao2024tau}와 같은 벤치마크의 핵심 관심사였다. 대부분의 노력은 예컨대 범용 LLM을 사용해 사용자 응답을 생성하거나 검증하는 방식~\cite{prabhakar2025apigen}으로 사용자 시뮬레이터에 감독 신호를 도입하는 데 초점을 맞추어 왔다. 반면 환경 자체를 사용해 사용자 시뮬레이터의 행동을 제약하고 형성하여 신뢰성을 높일 가능성에는 상대적으로 적은 관심이 주어졌는데, 이는 우리 접근의 핵심 원칙이다. 이러한 우려는 과제 지향 대화 시스템의 맥락에서 폭넓게 연구되어 왔으며, 초기 연구인 \cite{pietquin_survey_2013}은 사용자 시뮬레이션 평가 지표에 대한 포괄적 조사를 제공했다. 더 최근에는 \cite{kazi_large_2024}가 과제 지향 대화 시스템 평가를 위한 사용자-에이전트로 LLM을 효과적으로 사용할 수 있음을 보였고, 신중한 프롬프팅과 상태 추적이 더 신뢰할 수 있고 문맥을 인식하는 사용자 시뮬레이션으로 이어질 수 있음을 보여 주었다.

\paragraph{다중 에이전트 벤치마크.}
우리 연구는 다중 에이전트 프레임워크를 구축하고 평가하려는 노력~\cite{zhu2025multiagentbench}과도 관련이 있다.
우리 논문의 사용자와 에이전트를 다중 에이전트 시스템을 형성하는 것으로 볼 수 있지만, 우리 경우의 핵심 차이는 최종 평가가 여전히 사용자로부터 올바른 정보를 이끌어 내고 과제를 해결하기 위해 올바른 행동을 수행하는 에이전트의 능력에 초점을 둔다는 점이다. 이는 에이전트와 사용자 사이에 내재적 비대칭성을 도입한다. 우리의 초점은 순수한 다중 에이전트 문제를 해결하는 것이 아니라, 역시 행위 주체성을 가진 사용자를 효과적으로 안내하고 협업하는 에이전트의 역량에 있다. 이런 의미에서 이 프레임워크는 협력적(예: 문제 해결), 경쟁적(예: 구독 협상), 또는 하이브리드일 수 있으며, 에이전트가 사용자 실수나 오류까지 고려하여 시나리오를 적절히 식별하고 헤쳐 나갈 것을 요구한다.
